\subsubsection{A numerical condition forcing a formal cone}

\begin{lemma}[Proved under stated hypotheses: formal linearization of a curve]
\label{D:negative-curve-linearization}
Let $C$ be a smooth proper connected curve in a smooth surface
$Y$ over an algebraically closed field $k$, of any characteristic.
Write $I=\mathcal O_Y(-C)$ and
$N=N_{C/Y}=(I/I^2)^\vee$.  Suppose
\[
 \begin{aligned}
 H^1(C,T_C\otimes N^{-n})&=0 &&(n\geq1),\\
 H^1(C,N^{1-n})&=0 &&(n\geq2).
 \end{aligned}
 \tag{L}
\]
Then the formal completion of $Y$ along $C$ is isomorphic to the
formal completion of the total space of $N$ along its zero
section.  The isomorphism can be chosen to induce the identity
on $C$ and on its conormal bundle.  In particular, if $g$ is
the genus of $C$ and $C^2=-d<0$, the numerical condition
\[
                   d>\max\{0,4g-4\}
 \tag{N}
\]
implies this conclusion.
\end{lemma}

\begin{proof}
Put $C_n=(C,\mathcal O_Y/I^n)$ for $n\geq1$; thus $C_1=C$.
Let $V=\operatorname{Tot}(N)$, let $J$ be its zero-section ideal,
and write $V_n=(C,\mathcal O_V/J^n)$.  We construct compatible
isomorphisms $f_n:C_n\to V_n$ inducing the specified identities.
The distinction between the step $n=1$ and all later steps is
essential for the exact obstruction groups.

First consider the square-zero extension
\[
 0\longrightarrow N^{-1}\longrightarrow\mathcal O_{C_2}
   \longrightarrow\mathcal O_C\longrightarrow0.
\]
Local $k$-algebra sections of its last map exist because $C$ is
smooth.  The difference of two sections is a $k$-derivation
$\mathcal O_C\to N^{-1}$, and adding any such derivation to a
section gives another section.  Thus the sheaf of local sections
is a torsor under $T_C\otimes N^{-1}$.  These statements are the
infinitesimal lifting criterion and square-zero derivation
calculation in
\cite[Lemmas 37.11.7 and 17.28.11, tags 02H6 and 01UP]
     {StacksBNegativeCurve}.
The first vanishing in (L), at $n=1$, makes this torsor trivial.
A global section splits the extension as a square-zero
$\mathcal O_C$-algebra, giving
\[
 \mathcal O_{C_2}\simeq\mathcal O_C\oplus N^{-1}
                       =\mathcal O_{V_2}.
\]
Using the specified identification of $I/I^2$ with $N^{-1}$
gives $f_2$.  In particular, no vanishing of $H^1(C,\mathcal O_C)$
is required at this step.

Now let $n\geq2$ and suppose $f_n$ has been constructed.
Regard it as a morphism $C_n\to V$.  The immersion
$C_n\subset C_{n+1}$ has square-zero ideal
\[
 I^n/I^{n+1}=N^{-n},
\]
annihilated by the ideal of $C$ in $C_n$.  Since $V$ is smooth
over $k$, $f_n$ lifts Zariski locally to $C_{n+1}\to V$.
Differences of local lifts form the sheaf
\[
 \mathcal H_n
  =\mathcal Hom_{\mathcal O_{C_n}}
       (f_n^*\Omega_{V/k},I^n/I^{n+1})
  =(T_V|_C)\otimes N^{-n}.
\]
The sheaf of lifts is a torsor under $\mathcal H_n$;
see \cite[Lemmas 37.9.4--5 and 37.11.7,
tags 04FH, 04FJ and 02H6]{StacksBNegativeCurve}.
The projection $V\to C$ and its zero section give the canonical
splitting
\[
 T_V|_C\simeq T_C\oplus N.
\]
Consequently the exact obstruction group at this step is
\[
 H^1(C,\mathcal H_n)
  =H^1(C,T_C\otimes N^{-n})\oplus H^1(C,N^{1-n}),
                  \qquad n\geq2.
 \tag{O}
\]
It vanishes by (L), so a global lift $f_{n+1}:C_{n+1}\to V$
exists and agrees with the previously chosen $f_n$.

Because this lift restricts to the identity $C\to C$, it sends
the zero-section ideal $J$ into $I/I^{n+1}$.  It therefore
factors through $V_{n+1}$.  Its map on $J/J^2$ agrees with that
of $f_2$ and is the prescribed identity.  Both embeddings are
effective Cartier divisors, so multiplication identifies
\[
 J^r/J^{r+1}\simeq N^{-r}\simeq I^r/I^{r+1}.
\]
The resulting map on every associated graded piece, for
$0\leq r\leq n$, is an isomorphism.  Induction on the finite
ideal filtrations shows that $f_{n+1}:C_{n+1}\to V_{n+1}$ is
an isomorphism.  Passing to the inverse limit of these
compatible isomorphisms proves formal linearization.
All constructions used square-zero extensions and ordinary
derivations; no division by an integer is involved.

We finally verify (N).  Serre duality on the smooth proper
curve, with $\omega_C=\Omega_{C/k}$ and $\deg\omega_C=2g-2$,
gives
\[
 \begin{aligned}
 H^1(C,T_C\otimes N^{-n})^\vee
    &\simeq H^0(C,\omega_C^{\otimes2}\otimes N^n),
       &\deg(\omega_C^{\otimes2}\otimes N^n)&=4g-4-nd,\\
 H^1(C,N^{1-n})^\vee
    &\simeq H^0(C,\omega_C\otimes N^{n-1}),
       &\deg(\omega_C\otimes N^{n-1})&=2g-2-(n-1)d.
 \end{aligned}
\]
The exact primary inputs are
\cite[Lemmas 53.4.1 and 53.5.2, tags 0BS2 and 0BS6]
     {StacksBNegativeCurve}.
Under (N), the first degree is negative for every $n\geq1$,
and the second is negative for every $n\geq2$.
A line bundle of negative degree on a proper integral curve
has no nonzero section, proving (L).
\end{proof}

\begin{corollary}[Proved under stated hypotheses: sufficiently negative exceptional curves]
\label{D:negative-curve-tame}
Let $(S,s)$ be a normal surface germ over an algebraically closed
field of characteristic $p>0$.  Suppose a proper resolution of
this germ has exceptional fibre consisting of one smooth proper
curve $C$, and write $g=g(C)$ and $C^2=-d<0$.  If (L) holds for
$N_{C/Y}$, in particular if $d>\max\{0,4g-4\}$, every finite
normal cover of $S$ which is \'etale on $S\setminus\{s\}$ is
tame at the valuation defined by $C$.  This includes separability
of the residue-field extension and has no restriction on the
degree of the cover or on the Cartier index of $K_S$.
\end{corollary}

\begin{proof}
The resolution is smooth over the algebraically closed field
in a neighborhood of $C$.  Lemma~\ref{D:negative-curve-linearization}
identifies its formal neighborhood with the formal zero section
of the negative line bundle $N_{C/Y}$.  Apply
Proposition~\ref{D:graded-cone-tame}.  Its completion comparison
uses the genuine completed boundary DVR and proves the claimed
tameness for the original valuation.  We do not infer a Cartier
property of the canonical divisor and do not use an index-one
cover.
\end{proof}

\begin{remark}[Scope of the numerical condition]
For $g\geq2$ the numerical threshold is $d>4g-4$.
It is a sufficient condition, not a necessary condition for
linearization or tameness.  The cohomological criterion (L)
also covers borderline degrees when its indicated spaces of
sections vanish.  In particular, a failure of this numerical
inequality is not evidence of a wild cover.  Conversely, the
lemma does not assume in advance that the germ is a cone: it
constructs the required formal identification by compatible
infinitesimal lifts.
\end{remark}
